
% ----------------------------------------------------------------------
%                   LATEX TEMPLATE FOR PhD THESIS
% ----------------------------------------------------------------------

% based on Harish Bhanderi's PhD/MPhil template, then Uni Cambridge
% http://www-h.eng.cam.ac.uk/help/tpl/textprocessing/ThesisStyle/
% corrected and extended in 2007 by Jakob Suckale, then MPI-CBG PhD programme
% and made available through OpenWetWare.org - the free biology wiki


%: Style file for Latex
% Most style definitions are in the external file PhDthesisPSnPDF.
% In this template package, it can be found in ./Latex/Classes/
\documentclass[twoside,12pt]{Latex/Classes/PhDthesisPSnPDF}
\usepackage{acro}

%: Macro file for Latex
% Macros help you summarise frequently repeated Latex commands.
% Here, they are placed in an external file /Latex/Macros/MacroFile1.tex
% An macro that you may use frequently is the figuremacro (see introduction.tex)
% This file contains macros that can be called up from connected TeX files
% It helps to summarise repeated code, e.g. figure insertion (see below).

% insert a centered figure with caption and description
% parameters 1:filename, 2:title, 3:description and label
\newcommand{\figuremacro}[3]{
	\begin{figure}[htbp]
		\centering
		\includegraphics[width=1\textwidth]{#1}
		\caption[#2]{\textbf{#2} - #3}
		\label{#1}
	\end{figure}
}

% insert a centered figure with caption and description AND WIDTH
% parameters 1:filename, 2:title, 3:description and label, 4: textwidth
% textwidth 1 means as text, 0.5 means half the width of the text
\newcommand{\figuremacroW}[4]{
	\begin{figure}[htbp]
		\centering
		\includegraphics[width=#4\textwidth]{#1}
		\caption[#2]{\textbf{#2} - #3}
		\label{#1}
	\end{figure}
}

% inserts a figure with wrapped around text; only suitable for NARROW figs
% o is for outside on a double paged document; others: l, r, i(inside)
% text and figure will each be half of the document width
% note: long captions often crash with adjacent content; take care
% in general: above 2 macro produce more reliable layout
\newcommand{\figuremacroN}[3]{
	\begin{wrapfigure}{o}{0.5\textwidth}
		\centering
		\includegraphics[width=0.48\textwidth]{#1}
		\caption[#2]{{\small\textbf{#2} - #3}}
		\label{#1}
	\end{wrapfigure}
}

% predefined commands by Harish
\newcommand{\PdfPsText}[2]{
  \ifpdf
     #1
  \else
     #2
  \fi
}

\newcommand{\IncludeGraphicsH}[3]{
  \PdfPsText{\includegraphics[height=#2]{#1}}{\includegraphics[bb = #3, height=#2]{#1}}
}

\newcommand{\IncludeGraphicsW}[3]{
  \PdfPsText{\includegraphics[width=#2]{#1}}{\includegraphics[bb = #3, width=#2]{#1}}
}

\newcommand{\InsertFig}[3]{
  \begin{figure}[!htbp]
    \begin{center}
      \leavevmode
      #1
      \caption{#2}
      \label{#3}
    \end{center}
  \end{figure}
}


%%% Local Variables: 
%%% mode: latex
%%% TeX-master: "~/Documents/LaTeX/CUEDThesisPSnPDF/thesis"
%%% End: 

\usepackage{nicefrac}
\usepackage[inner=0.75in,outer=0.65in,top=0.75in,bottom=0.75in]{geometry}
\usepackage[sc]{mathpazo}
\usepackage{multirow}


%: ----------------------------------------------------------------------
%:                  TITLE PAGE: name, degree,..
% ----------------------------------------------------------------------
% below is to generate the title page with crest and author name

%if output to PDF then put the following in PDF header
\ifpdf  
    \pdfinfo { /Title  (PhD and MPhil Thesis Classes)
               /Creator (TeX)
               /Producer (pdfTeX)
               /Author (YourName your@email.net)
               /CreationDate (D:YYYYMMDDhhmmss)  %format D:YYYYMMDDhhmmss
               /ModDate (D:YYYYMMDDhhmm)
               /Subject (xyz)
               /Keywords (add, your, keywords, here) }
    \pdfcatalog { /PageMode (/UseOutlines)
                  /OpenAction (fitbh)  }
\fi


\title{Title of your thesis}



% ----------------------------------------------------------------------
% The section below defines www links/email for author and institutions
% They will appear on the title page of the PDF and can be clicked
\ifpdf
  \author{\href{mailto:your@email.net}{YourName}}
%  \cityofbirth{born in XYZ} % uncomment this if your university requires this
%  % If city of birth is required, also uncomment 2 sections in PhDthesisPSnPDF
%  % Just search for the "city" and you'll find them.
  \collegeordept{\href{http://www.something.net}{CollegeOrDepartment}}
  \university{\href{http://www.something.net}{University}}

  % The crest is a graphics file of the logo of your research institution.
  % Place it in ./0_frontmatter/figures and specify the width
  \crest{\includegraphics[width=4cm]{logo}}
  
% If you are not creating a PDF then use the following. The default is PDF.
\else
  \author{YourName}
%  \cityofbirth{born in XYZ}
  \collegeordept{CollegeOrDept}
  \university{University}
  \crest{\includegraphics[width=4cm]{logo}}
\fi

%\renewcommand{\submittedtext}{change the default text here if needed}
\degree{Bachelor of Science }
\degreedate{year month}


% ----------------------------------------------------------------------
       
% turn of those nasty overfull and underfull hboxes
\hbadness=10000
\hfuzz=50pt


%: --------------------------------------------------------------
%:                  FRONT MATTER: dedications, abstract,..
% --------------------------------------------------------------

\begin{document}

%\language{english}

% sets line spacing
\renewcommand\baselinestretch{1.2}
\baselineskip=18pt plus1pt


%: ----------------------- generate cover page ------------------------

\maketitle  % command to print the title page with above variables


%: ----------------------- cover page back side ------------------------
% Your research institution may require reviewer names, etc.
% This cover back side is required by Dresden Med Fac; uncomment if needed.



%: ----------------------- abstract ------------------------

% Your institution may have specific regulations if you need an abstract and where it is to be placed in the document. The default here is just after title.


% Thesis Abstract -----------------------------------------------------

%\begin{abstractslong}    %uncommenting this line, gives a different abstract heading
\begin{abstracts}        %this creates the heading for the abstract page
This project is to make it easier for travelers to find the required information on transport terminals in Ghana for effective trip planning and decision making. 
\end{abstracts}
%\end{abstractlongs}


% ---------------------------------------------------------------------- 


% The original template provides and abstractseparate environment, if your institution requires them to be separate. I think it's easier to print the abstract from the complete thesis by restricting printing to the relevant page.
% \begin{abstractseparate}
%   
% Thesis Abstract -----------------------------------------------------

%\begin{abstractslong}    %uncommenting this line, gives a different abstract heading
\begin{abstracts}        %this creates the heading for the abstract page
\addcontentsline{toc}{chapter}{ABSTRACT}
Information about public transport operators and terminals in Ghana and most Africa countries, where transportation by road is a major means in relation to routes they operate, departure times, fares and location is not easy to come by when planning trip ahead unless you visit these terminals where information is readily available in most cases; sometimes relying on word of mouth from other travelers which could be valid or invalid. This project seeks to make it easier for travelers to find the required information on transport terminals in Ghana for effective trip planning and decision making. In conclusion a public transport search web portal application was developed. The application was developed using data collected and mapped into OpenStreetMap database, the data was extracted, cleaned and analysed in QGIS and used to create a geodatabase; using PostgreSQL and PostGIS to suit the needs of this project. The geodatabase can also be accessed by other stakeholders externally without interacting with the web application.
\end{abstracts}
%\end{abstractlongs}


% ---------------------------------------------------------------------- 

% \end{abstractseparate}


%: ----------------------- tie in front matter ------------------------

\frontmatter
% Thesis Dedictation ---------------------------------------------------

\begin{dedication} %this creates the heading for the dedication page
This project wok is dedicated to my mother, Prince Ahiabu, Jemila Adams and Okatachie Afrifa Amankwaa. Your unending love and supports is highly appreciated.  

\end{dedication}

% ----------------------------------------------------------------------

% Thesis Acknowledgements ------------------------------------------------


%\begin{acknowledgementslong} %uncommenting this line, gives a different acknowledgements heading
\begin{acknowledgements}      %this creates the heading for the acknowlegments
I am grateful to Dr Hamidu Abdel-Fatao, my project supervisor for guiding me through the writing of this report, and my lecturers who equipped me with the knowledge required to write it.
\end{acknowledgements}
%\end{acknowledgmentslong}

% ------------------------------------------------------------------------





%: ----------------------- contents ------------------------

\setcounter{secnumdepth}{3} % organisational level that receives a numbers
\setcounter{tocdepth}{3}    % print table of contents for level 3
\tableofcontents            % print the table of contents
% levels are: 0 - chapter, 1 - section, 2 - subsection, 3 - subsection


%: ----------------------- list of figures/tables ------------------------

\listoffigures	% print list of figures

\listoftables  % print list of tables


%: ----------------------- glossary ------------------------

% Tie in external source file for definitions: /0_frontmatter/glossary.tex
% Glossary entries can also be defined in the main text. See glossary.tex
%% this file is called up by thesis.tex
% content in this file will be fed into the main document

% Glossary entries are defined with the command \nomenclature{1}{2}
% 1 = Entry name, e.g. abbreviation; 2 = Explanation
% You can place all explanations in this separate file or declare them in the middle of the text. Either way they will be collected in the glossary.

% required to print nomenclature name to page header
\markboth{\MakeUppercase{\nomname}}{\MakeUppercase{\nomname}}


% ----------------------- contents from here ------------------------

% chemicals
\nomenclature{DAPI}{4',6-diamidino-2-phenylindole; a fluorescent stain that binds strongly to DNA and serves to marks the nucleus in fluorescence microscopy} 
\nomenclature{DEPC}{diethyl-pyro-carbonate; used to remove RNA-degrading enzymes (RNAases) from water and laboratory utensils}
\nomenclature{DMSO}{dimethyl sulfoxide; organic solvent, readily passes through skin, cryoprotectant in cell culture}
\nomenclature{EDTA}{Ethylene-diamine-tetraacetic acid; a chelating (two-pronged) molecule used to sequester most divalent (or trivalent) metal ions, such as calcium (Ca$^{2+}$) and magnesium (Mg$^{2+}$), copper (Cu$^{2+}$), or iron (Fe$^{2+}$ / Fe$^{3+}$)}



 

\begin{multicols}{2} % \begin{multicols}{#columns}[header text][space]
\begin{footnotesize} % scriptsize(7) < footnotesize(8) < small (9) < normal (10)

\printnomenclature[1.5cm] % [] = distance between entry and description
\label{nom} % target name for links to glossary

\end{footnotesize}
\end{multicols}



%: --------------------------------------------------------------
%:                  MAIN DOCUMENT SECTION
% --------------------------------------------------------------

% the main text starts here with the introduction, 1st chapter,...
\mainmatter

\renewcommand{\chaptername}{} % uncomment to print only "1" not "Chapter 1"


%: ----------------------- subdocuments ------------------------

% Parts of the thesis are included below. Rename the files as required.
% But take care that the paths match. You can also change the order of appearance by moving the include commands.


% this file is called up by thesis.tex
% content in this file will be fed into the main document

%: ----------------------- introduction file header -----------------------
\chapter{Introduction}

% the code below specifies where the figures are stored
\ifpdf
    \graphicspath{{1_introduction/figures/PNG/}{1_introduction/figures/PDF/}{1_introduction/figures/}}
\else
    \graphicspath{{1_introduction/figures/EPS/}{1_introduction/figures/}}
\fi

% ----------------------------------------------------------------------
%: ----------------------- introduction content ----------------------- 
% ----------------------------------------------------------------------



%: ----------------------- HELP: latex document organisation
% the commands below help you to subdivide and organise your thesis
%    \chapter{}       = level 1, top level
%    \section{}       = level 2
%    \subsection{}    = level 3
%    \subsubsection{} = level 4
% note that everything after the percentage sign is hidden from output

\section{Problem Statement}
Road transport service provided by both formal and informal sectors in Ghana is by far the most popular and principal means of conveying passengers, goods and other services between any two or more locations \citep{aidoo_passengers_2013}. Some report have it that, road transport carries over $ 95\% $ of all passenger and freight traffic and about $ 97\% $ of all passenger miles in Ghana \citep[p.~195]{unesco_transportation:_????}, a country that is experiencing rapid demographic and economic growth. The vast majority of passengers commuting between places, be it \textit{intra-city} or \textit{inter-city}, mostly rely on public transport services in the form of privately owned or corporate taxis, \textit{tro tros} (shared minivans), buses commuting between major cities.

In spite of the heavy reliance on public road transport services by the general populace in Ghana, finding transport terminals which offer reliable road transport services is not as easy as it should be. The difficulty in finding transport terminals is attributable to the fact that little or no information about the availability of transport services and their locations is accessible to the public. Additionally, the non-existent of a means to compare transport fares by various service providers often makes it difficult for the potential passenger to make the right choices. It is the goal of every potential passenger to find the fastest, safest and most cost-efficient means of transiting from one location to another. 

Inspired by the aforementioned shortcomings of the existing public road transportation system, this project seeks to develop road transport terminal search tool
aimed at mitigating, if not eliminate entirely, these problems with the public road transport industry.

\newpage
\section{Project Objectives}
The objective of this project include the following:
\begin{itemize}
	\item To develop a web application that provides detailed information about transport terminals in Ghana to mitigate the difficulty in finding transport terminals and also to provide a platform for to compare other factors by travelers.
	%	\item To provide \hl{reusable data} that can be accessed in mobile and desktop applications that support Geo Uniform Resource Identifier (URI) scheme.
	%	\item With Global Positioning System (GPS) supported devices, the work also provides a means of navigating to destination based by incorporating existing mapping solution such as
	%    \begin{itemize}[label=$\circ$]
	%	\item OSMAnd
	%	\item MAPS.ME
	%	\item Google Maps
	%	\item Apple Maps
	%	\item Marble
	%	\item Google Earth
	%	\end{itemize}
\end{itemize}

\section{Project Outcomes}
The following will be achieved at the end of the project:
\begin{itemize}
	\item An web application;
%	\item \sout{Possibly generate high definition printer friendly maps that serves nearly same purpose as the web platform.}
	\item A detail map of selected transport terminals on OpenStreetMap.
\end{itemize}


\section{Methods to be used}
The methods to be used for the project are as below:
\begin{itemize}
	\item Literature review on proposed topic
	\item Study and understanding of online maps (creating, updating, deleting);
	\item The system will be developed on a handful of local machines, but with scalability and ease of deployment on any kind of infrastructure;
	\item The system will be prototyped in Python programming language and Django web framework. If this language proves good enough for deployment purposes it will be used in the final product;
	\item Survey and crowd-sourcing information on some transport terminals to facilitate database creation;
	\item QGIS will be used to clean and analyze data collected;
	\item Spiral software development model.
\end{itemize}

\section{Tools and Facilities to be used}
The facilities required for this project include:
\begin{itemize}
	\item University of Mines and Technology (UMaT) library;
	\item Internet;
	\item General search engines as such Wikipedia and DuckDuckGo;
	\item Open Source software repositories such as GitHub;
	\item OpenStreetMap;
	\item Documentation of any software or libraries used.
\end{itemize}

\section{Scope of Work}
This works seeks to aid travelers acquire detailed information on selected bus terminals in Ghana. This is to enable better trip planning, reduce time taken finding these terminals and also allowing passengers to compare transport fares and choosing the best option they can afford.
Further more the application only indicates source and destination terminals and works in the any modern web browser such as Mozilla Firefox or Google Chrome.

\section{Organization of Project}
This project is divided into five chapters. The first chapter talks about the problem to be solved, the objectives, the methods, tools and facilities used, and project outcomes. The second chapter discusses relevant literature and related works. The third chapter discusses how the problem was solved. The fourth chapter talks about the operation of the developed application. The project concludes in the fifth chapter where, limitations and recommendations for future improvements were discussed.

	% background information
% this file is called up by thesis.tex
% content in this file will be fed into the main document

\chapter{Literature Review} % top level followed by section, subsection


% ----------------------- contents from here ------------------------

\section{Background Study}
Transportation is a system or means of transporting people or goods from one place to another. Transportation is key to the movement of people, goods and services leading to development. In Sub-Sahara and most African countries, the major means of transportation is by road transport.

Transportation in Ghana dates back to precolonial era. Most major roads and railway lines were constructed by the Europeans to facilitate the movement of natural resources and raw materials such as timber, gold and bauxite to generating communities to the coastal areas. Transport in all its modes is one of the most important catalysts for development \textit{(cite)}. %http://www.mrh.gov.gh/files/publications/National_Transport_Policy___December__2008.pdf
After independence road transportation remains a major player in Ghana's economy connecting all ten administrative regions of Ghana. 

\section{Categories of Road Transport}
Road Transport in Ghana can be grouped into 4 main categories:\textit{(cite)}
\begin{itemize}
	\item Urban
	\item Express Services
	\item Rural - Urban
	\item Rural
\end{itemize}

\subsection{Urban}
Urban is used mostly by residents of urban areas commuting to and from work, school and other places of convenience. It is usually carried out by private taxis, personal cars, \textit{trotros} and state owned transport services.

\subsection{Express Services}
Express services are demanded services by travelers; this could be in the form of taxis, buses and minivans.

\subsection{Rural-Urban}
Rural-Urban transportation is by far the most popular group of public transportation in Ghana as well as a major factor in Rural-Urban migration. Market trading in the urban centres of Ghana is a predominantly female economic activity and is a fundamental element in the survival strategies of many low-income households. Petty trading is the predominant form of commercial activity and as such, given the financial constraints inherent in this form of trading, necessitates regular travel between wholesale markets and selling places on the part of female traders. Female traders make use of the public transport system, in combination with supplementary services such as portering, to meet their travel needs \citep{grieco1995informal}. 

\subsection{Rural}
This is transportation mainly within rural areas. Depending on the status of each rural area transportation could range for donkey pulled carts, foot, bicycles and motorcycles, taxis and minivans in average conditions.

\subsection{Owners or Operators}
Transportation service operators or owners can be classified as:
\begin{itemize}
	\item Private
	\item Government
	\item Private-Government
\end{itemize}

\section{The Industry Structure}
Within the main metropolitan areas of Ghana, there are two main forms of public transport operations \citep{finn2009new}:
\begin{enumerate}
	\item \textit{Tro-tro} (mini buses) and shared taxi services, which are managed by unions and cooperatives and offer services along defined routes, usually between terminals or ‘lorry parks/stations’. These operations suffer from a number of quality problems including:
	
	\begin{enumerate}
		\item operation of a ‘fill and go’ system which can result in long delays for users in the offpeak, and difficulty to board along the route;
		\item large numbers of vehicles parked at
		terminals in the off-peak leading to congestion, inefficiency, and long hours for
		drivers; 
		\item lack of incentives for vehicle owners to improve their vehicles or to train their drivers properly.
	\end{enumerate}
	
	\item Large bus services, mostly provided by the new Metro Mass Transit (MMT), a quasi-private company that receives favorable financial support from the government.
\end{enumerate}

\subsection{Operating unions and associations}. 
Individually or privately operated transport services are members of unions or associations. These unions and associations serve as regulatory and mouth-piece to each of their members \citep{fouracre1994public}. There are three major operating unions and associations.

\subsection{Ghana Private Road Transport Union}
The Ghana Private Road Transport Union (GPRTU), a national union, is reported to have about 90\% of the \textit{tro-tro} and shared taxi business. The fundamental units are Locals, which operate the individual routes, and Branches, which are regional clusters of Locals. GPRTU represents the interests of both drivers and of vehicle owners.

\subsection{Progressive Transport Owners Association}
The Progressive Transport Owners Association (PROTOA), is a national association that operates both tro-tro and shared taxi business and is organized along the same structure as the GPRTU. PROTOA mainly represents the interest of owners. 

\subsection{Ghana Co-operative Transport Association}
The Ghana Co-operative Transport Association (GCTA) is a national association also organized along the lines of GPRTU and represents interest of both owners and drivers.

\subsection{Other private operators}
Other private operators, such as Agate, Kingdom Transport, and Pergah Transport, are companies operating several buses and offering a range of services including contract service, urban services, and intercity services.

\subsection{Ghana Road Transport Coordination Council}
Ghana Road Transport Coordination Council (GRTCC) is an umbrella body of all transport operators in Ghana, including the unions and associations, other locally based associations, and other operators (both passenger and road haulers). GRTCC represents the interests of road transport operators, especially in negotiating with the Government of Ghana for transport tariffs and assistance in acquisition of buses.

\section{Operators}
Transport Operators are individual, state owned or both private and government partnerships managing the affairs of a particular brand.

\subsection{Intercity STC Limited}
Intercity bus transport is a popular means of travelling between cities and aligning villages and towns in Ghana. Its services include freight and passenger movements from one location to the other. For this service to be provided, a company has to be formed. As a result of that there have been concerted attempts by various past Governments of Ghana to offer intercity bus transport service to her citizen. One of such efforts is the establishment of Intercity State Transport Company (ISTC). But there is a number of private transport operators 
of which Ghana Private Road Transport Union (GPRTU) offers about $70-80\%$ of passenger and freight traffic. This is an off shoot of intra urban dominance of GPRTU of $70-80\%$ \citep{abane2011travel}. GPRTU has been able to co-opt other intercity bus transport operators by sharing some of its terminals/stations with other transport 
companies/union such as VVIP/VIP, DIPLOMAT. Aside this, some private owners or operators like VIP/VVIP, DIPLOMAT for instance are either members or former members/executives of the union. Other unions/transport operators in the industry are Concerned Drivers Union, Progressive Transport Owners Association and Co-operative \citep{ojobus}.

\textit{Insert map of STC Terminals across Ghana}

\subsection{Metro Mass Transit}
Metro  Mass  Transit  Limited, Ghana was established in 2001 by the former President of Ghana, John Kuffour who directed the re-introduction of public mass transport in the metropolitan and municipal areas to ensure safe, affordable, efficient and reliable movement of Ghanaians. Since then, the Government has been actively promoting public mass transportation \citep{olateju2009appraisal}. MMT receives financial support from the Government and currently operates about 500 buses of which some 200 operate in the greater Accra area \citep{finn2009new}.

\subsection{Aayalolo}
Aayalolo Bus Rapid Transit(BRT) system has recently been incorporated. Since November 2016, the company has been running three services on the Amasaman-Tudu/
Accra Central corridor. There are plans to roll out the Adenta-Tudu/Accra Central corridor and subsequently, another service along the Tema Beach Road-Tudu/Accra Central corridor. There are plans to replicate the BRT mass transit services in other major Ghanaian cities. Prior to the Aayalolo bus service \citep{agyemang2017mode}.

\section{Maps}
A map is a graphic representation or scale model of spatial concepts. It is a means for conveying geographic information. Maps are a universal medium for communication, easily appreciated and understood by most people, regardless of language or culture. Maps record the geographical information that is fundamental to reconstructing past places, towns, even cities.

\section{Digital mapping}
Digital mapping is the process by which a collection of data is compiled and converted into a virtual image. The primary function of this technology is to produce maps that give accurate representation of a particular area, detailing major road arteries and other points of interest. The technology also allows the calculation of distances from one place to another.
The roots of digital mapping lie within traditional paper maps. Paper maps provide basic landscapes similar to digitized road maps, yet are often cumbersome, cover only a designated area, lack many specific details such as road blocks etc. In addition, there is no way to update a paper map except to obtain a new version. Conversely, digital maps, in many cases, can be updated.
Early digital maps had the same functionality as traditional maps-that is, they provided a ‘virtual view’ of roads generally outlined by the terrain encompassing the mounding area. However, as digital maps have grown with the expansion of GPS technology in the past decade, live traffic updates, points of interests and service locations have been added to enhance digital maps to be more ‘user conscious’ \textit{cite}. Digital maps heavily rely on a vast amount of data collection over time.

\section{Existing Technologies}
Digital maps have changed the perception of maps and introduced much flexibility compared to paper maps. Existing technologies such as OpenStreetMap, Google Maps, Bing Maps and Taximap provides web map services and that captures transportation information but not into much detail. Taximap on the other seeks to localize this \citep{vinet2014arch}.

\subsection{OpenStreetMap}
OpenstreetMap (OSM) is a collaborative project started in England in 2004 by Steve Coast. The aim of OSM is to create and provide free geographic data. The project aims to compensate the lack of free data because geographic data, even freely available, are provided with licenses restricting the use of information and the creativity according to project leaders. The data are distributed under the license “Creative Commons Attribution-ShareAlike 2.0 license”. This license allows using the data completely freely, in condition to distribute any derived data under the same license. For instance, corrected OSM data cannot be sold.
Data stored in OSM by contributors of the project are modelled and stored in tagged geometric primitives. For example, a road is a polyline with tags \textit{highway= “primary”, oneway= “no” and name= “N10”}. Geometric primitives are of three types: points, paths (polylines) and relationships (linking points and paths with tags) that are not really geometric primitives. The surfaces are represented by close paths. Data are available from any area specified for export in a specific XML based format. It has to be translated if anyone wishes to use the data in another application. Data is captured using GIS software adapted to OSM data with editing functions to create OSM geometric primitives and tag them. Different software exists to edit and capture OSM data (Potlatch, JOSM, Merkaartor). OSM applications currently aim to foster mapping creativity of potential contributors of geographic data \citep{girres_quality_2010}.

\subsection{Google Maps}
Google Maps is a propitiatory tool for navigation. Google Maps is used by many people around the world. Google has an online database of structures; which is better covered in most developed countries compared to developing countries. Searching for places of interests is done by geocoding - \textit{converting names or addresses of places to a location on a map} or reverse geocoding - \textit{converting latitudes and longitudes to a readable address or name}. Google Maps is made possible through Volunteered Geographic Information (VGI) and other third party proprietary data sources. Adding, modifying and deleting features from Google Maps has recently become difficult for a new contributor. Google Maps like OpenStreetMap could also fail to pin point the exact location of a terminal.

% Clound route from one location to another but relative to position of towns not the various bus terminals

\subsection{Taximap}
Taximap is a social enterprise start-up and public transport search portal in South Africa. It provides information about minibus taxi routes, fares and operating hours with commuters. Information about on over 800 mapped routes are available. The platform rely on user feedback to keep the routes up-to-date and accurate, by encouraging users to leave feedbacks. Taximaps does well by telling the exact Taximap has one drawback by not allowing users of the platform to specify both their departure and destination locations. 


\section{Proposed System}

\section{Software Review}

\subsection{Arch Linux}
The Arch Linux operating system was most suitable for my project. Arch Linux is a native Linux based operating system produced for computers with chipsets based on i686 and x86-64 architectures (Griffin, 2002); it was chosen because of its simplicity, community involvement and a well documented wiki and community support \citep{vinet2014arch}.

\subsection{Python}
Python is an object oriented, interpreted programming language. It runs on a wide variety of systems which include Linux, Unix, Mac OS, Windows, BSD, etc. It also has several implementations, such as IronPython which runs on .NET CLR, Jython, amongst others \citep{van2007python}. 

\subsection{Django}
Django is a popular Python Web Development Framework. It is dubbed the web develoment framework from developers with deadline. Django was born from the newsroom as a framework for development new and media related web applications.

\subsection{Emacs}
Emacs was the development environment used to develop this project. It was initially released in 1976 and active development continues to date. Its flexibility is one of its major advantages. Amongst other things, it provides a text editor with efficient key bindings and a method for interacting with inferior shells such as Python shell. What this means is that, after a line code has been written, it can be directly sent to the Python interpreter, where the code is executed and the results are shown immediately. This makes exploratory development much more efficient and streamlined. 

\subsection{Java OpenStreetMap}
Java OpenStreetMap (JOSM) is an extensible editor for OpenStreetMap. It supports loading GPX tracks, background imagery and OSM data from local sources as well as from online sources and allows to edit the OSM data (nodes, ways, and relations) and their metadata tags. JOSM is an open source. JOSM was used for the initial processing of data collected and uploading into OpenStreetMap.

\subsection{PostgresSQL}
The database management system which is more suitable for my proposed system is PostgreSQL. PostgreSQL is an open source relational database management system that began as a University of California, Berkeley project. PostgreSQL was selected above MySQL and MSSQL (Microsoft SQL Server) because it's open source and great support for extensions. PostGIS extension provides a great support for geographic data. It also powers OpenStreetMap and Skype databases.

It has enterprise class features such as SQL windowing functions, the ability to create aggregate functions and also utilize them in window constructs, common table and
recursive common table expressions, and streaming replication. These features are rarely found in other open source database platforms, but commonly found in newer versions of the proprietary databases such as Oracle, SQL Server, and IBM DB2. What sets it apart from other databases, including the proprietary ones we just mentioned,
is the ease with which you can extend it without changing the underlying base—and in many cases, without any code compilation. Not only does it have advanced features,
but it performs them quickly. It can outperform many other databases, including proprietary ones for many types of database workloads.




% ---------------------------------------------------------------------------
% ----------------------- end of thesis sub-document ------------------------
% ---------------------------------------------------------------------------
						% aims of the project

`% this file is called up by thesis.tex
% content in this file will be fed into the main document

\chapter{Discussion} % top level followed by section, subsection


% ----------------------- paths to graphics ------------------------

% change according to folder and file names
\ifpdf
    \graphicspath{{7/figures/PNG/}{7/figures/PDF/}{7/figures/}}
\else
    \graphicspath{{7/figures/EPS/}{7/figures/}}
\fi


% ----------------------- contents from here ------------------------






% ---------------------------------------------------------------------------
% ----------------------- end of thesis sub-document ------------------------
% ---------------------------------------------------------------------------               % discussion of results


% this file is called up by thesis.tex
% content in this file will be fed into the main document

\chapter{Materials \& methods} % top level followed by section, subsection


% ----------------------- paths to graphics ------------------------

% change according to folder and file names
\ifpdf
    \graphicspath{{8/figures/PNG/}{8/figures/PDF/}{8/figures/}}
\else
    \graphicspath{{8/figures/EPS/}{8/figures/}}
\fi

% ----------------------- contents from here ------------------------




 

% ---------------------------------------------------------------------------
%: ----------------------- end of thesis sub-document ------------------------
% ---------------------------------------------------------------------------



 






        % description of lab methods




% --------------------------------------------------------------
%:                  BACK MATTER: appendices, refs,..
% --------------------------------------------------------------

% the back matter: appendix and references close the thesis


%: ----------------------- bibliography ------------------------

% The section below defines how references are listed and formatted
% The default below is 2 columns, small font, complete author names.
% Entries are also linked back to the page number in the text and to external URL if provided in the BibTex file.

% PhDbiblio-url2 = names small caps, title bold & hyperlinked, link to page 
\begin{multicols}{2} % \begin{multicols}{ # columns}[ header text][ space]
\begin{tiny} % tiny(5) < scriptsize(7) < footnotesize(8) < small (9)

\bibliographystyle{Latex/Classes/PhDbiblio-url2} % Title is link if provided
\renewcommand{\bibname}{References} % changes the header; default: Bibliography

\bibliography{9_backmatter/references} % adjust this to fit your BibTex file

\end{tiny}
\end{multicols}

% --------------------------------------------------------------
% Various bibliography styles exit. Replace above style as desired.

% in-text refs: (1) (1; 2)
% ref list: alphabetical; author(s) in small caps; initials last name; page(s)
%\bibliographystyle{Latex/Classes/PhDbiblio-case} % title forced lower case
%\bibliographystyle{Latex/Classes/PhDbiblio-bold} % title as in bibtex but bold
%\bibliographystyle{Latex/Classes/PhDbiblio-url} % bold + www link if provided

%\bibliographystyle{Latex/Classes/jmb} % calls style file jmb.bst
% in-text refs: author (year) without brackets
% ref list: alphabetical; author(s) in normal font; last name, initials; page(s)

%\bibliographystyle{plainnat} % calls style file plainnat.bst
% in-text refs: author (year) without brackets
% (this works with package natbib)


% --------------------------------------------------------------

% according to Dresden med fac summary has to be at the end
%
% Thesis Abstract -----------------------------------------------------

%\begin{abstractslong}    %uncommenting this line, gives a different abstract heading
\begin{abstracts}        %this creates the heading for the abstract page
This project is to make it easier for travelers to find the required information on transport terminals in Ghana for effective trip planning and decision making. 
\end{abstracts}
%\end{abstractlongs}


% ---------------------------------------------------------------------- 


%: Declaration of originality




\end{document}
